\documentclass[12pt,a4paper,psamsfonts]{amsart}
\usepackage{amssymb,mathtools}
\usepackage{calc}
\setlength{\topmargin}{0cm}
\setlength{\oddsidemargin}{0cm}
\setlength{\evensidemargin}{0cm}
\setlength{\marginparwidth}{0cm}
\setlength{\marginparsep}{0cm}
\setlength{\textheight}{\paperheight - 2in - 35pt}
\setlength{\textwidth}{\paperwidth - 2in}
\setlength{\headheight}{12.5pt}
\setlength{\headsep}{25pt}
\setlength{\footskip}{30pt}
\renewcommand{\baselinestretch}{1.3}
\pagestyle{headings}
\raggedbottom
\setlength{\emergencystretch}{2em}
\usepackage[colorlinks=true,urlcolor=blue,linkcolor=blue,citecolor=blue]{hyperref}
\newtheorem*{theorem}{Theorem}
\theoremstyle{definition}
\newtheorem*{setup}{Setup}
\theoremstyle{plain}
\newtheorem{lemma}{Lemma}[section]
\newtheorem{proposition}[lemma]{Proposition}
\newtheorem{corollary}[lemma]{Corollary}
\newcommand{\Q}{\mathbb Q}
\newcommand{\C}{\mathbb C}
\newcommand{\OO}{\mathcal O}
\newcommand{\mult}{\operatorname{mult}}
\newcommand{\nef}{\operatorname{n}}
\title{Nef Abundance: Framework -- Jet Separation}
\date{Temporary draft -- October 11, 2026}
\begin{document}
\maketitle

\section{Jet Separation}

Let $X$ be a smooth projective threefold over $\C$ such that
$K_X\sim0$ and $q(X)=0$. Let $D$ be a nef Cartier divisor on $X$
such that $\nef(X,D)=3$ and $D^3=0$.

\begin{setup}
\leavevmode
\begin{itemize}
\setlength{\itemsep}{2pt}
\setlength{\parsep}{0pt}
\item Fix $\epsilon\in\Q_{>0}$ sufficiently small such that
\[
 343\epsilon^3<\tfrac14,\qquad40\epsilon<\tfrac38.
\]
\item Fix an ample Cartier divisor $A$.
For integers $n\ge1$, put
\[
 a_n=((D+A/n)^3)^{1/3},\qquad
 r_n=\max\{1,\lfloor\epsilon/a_n+1/2\rfloor\},\qquad
 P_n=r_n(D+A/n).
\]
Then $P_n$ is ample and rational Cartier, $r_n\to\infty$, and
$P_n^3\to\epsilon^3$. Indeed, $a_n\to0$ and
$|r_na_n-\epsilon|\le a_n/2$ for all sufficiently large $n$.
\item Let $p_i:X\times X\to X$ be the projections, for $i=1,2$.
Put
\[
 Z=\mathbb P_{X\times X}
       \bigl(p_1^*\OO_X(D)\oplus p_2^*\OO_X(D)\bigr).
\]
Let $\pi:Z\to X\times X$ be the projection.
Let $S_1,S_2\subset Z$ be the two summand sections and set
\[
 U=Z\setminus(S_1\cup S_2).
\]
\item Set $\xi=c_1(\OO_Z(1))$. For a positive integer $s$, put
\[
 R_{n,s}=(p_1\circ\pi)^*P_n+(p_2\circ\pi)^*P_n+s\xi.
\]
Whenever $|kR_{n,s}|$ is used, $k>0$ is an integer with
$kR_{n,s}$ Cartier, equivalently $kP_n$ Cartier.
\end{itemize}

\end{setup}

\medskip
\noindent\textbf{Jet separation.}
Let $W$ be a smooth variety, $L$ a Cartier divisor, and $z\in W$.
For an integer $j\ge0$, $|L|$ separates $j$-jets at $z$ if
\[
 \forall u\in\OO_{W,z}(L)\quad
 \exists\sigma\in H^0(W,\OO_W(L))\quad
 \sigma_z-u\in\mathfrak m_z^{j+1}\OO_{W,z}(L).
\]
Here $\mathfrak m_z\subset\OO_{W,z}$ is the maximal ideal.

\begin{lemma}\label{lem:full-jets}
There exists a positive integer $s$ such that, for every $n\gg0$,
$|kR_{n,s}|$ separates $8k$-jets at some $z\in U$ for some
integer $k>0$ with $kR_{n,s}$ Cartier.
\end{lemma}

\medskip
\noindent\textbf{Proof outline for Lemma~\ref{lem:full-jets}.}
The proof has five steps.
\begin{enumerate}
\item Normalize $P_n$ so that $r_n\to\infty$ while $P_n^3$ stays
bounded and bounded away from zero. The projective-bundle calculation
in Lemma~\ref{lem:bundle-volume} then controls $R_{n,s}^7$.
\item Suppose no point of $U$ has Seshadri constant greater than $8$.
The moving-kernel construction produces base components $V$ with
high multiplicity. Lemma~\ref{lem:hd-base-degree} bounds their degrees;
finite subadjunction and Lemma~\ref{lem:canonical-centre} bound their
canonical intersections. These are assembled in
Lemma~\ref{lem:sevenfold-centres}, including the corresponding slice bounds.
\item Apply the lower-dimensional induction consequence
(Lemma~\ref{lem:lower-dimensional}) and the curve bound
(Lemma~\ref{lem:bounded-curves}) to exclude positive-dimensional
projection fibres. The remaining components are threefold
correspondences, generically finite over both factors of $X\times X$
(Lemma~\ref{lem:finite-correspondences}).
\item Exclude moving branch divisors and obtain finitely many Stein
covers. The graph argument reduces the remaining family to an abelian
model. Lemma~\ref{lem:bounded-correspondence-exclusion} then bounds
$r_n$ independently of $n$, contradicting $r_n\to\infty$.
\item Hence $\varepsilon(R_{n,s};z)>8$ for some $z\in U$.
Serre vanishing on the blowup at $z$ gives separation of $8k$-jets
for sufficiently divisible $k$.
\end{enumerate}

\medskip
\noindent\textbf{Draft status.}
The proof below is extracted from the existing threefold proof draft.
Its lower-dimensional hypotheses, model inputs, and source-dependent
arguments have not been independently verified. This note isolates the indicated framework lemma for discussion
and checking.

\section{Lower-dimensional induction hypotheses}

All varieties are over $\C$. In dimensions strictly less than three,
we assume the following two abundance statements for projective klt
$\Q$-pairs $(Y,B)$ and nef rational Cartier divisors $M$:
\begin{itemize}
\item If $K_Y+B$ is nef, then it is semiample.
\item If $K_Y+B$ is pseudo-effective and $K_Y+B+M$ is nef,
then $K_Y+B+M$ is numerically equivalent to a semiample
rational Cartier divisor.
\end{itemize}
These are the log abundance and nef abundance induction hypotheses.
Birational reductions use the terminating adjoint programs and effective
comparisons of the source framework \cite[Sections~2 and~8]{NS}.
They are model inputs, separate from the two abundance hypotheses.


\begin{lemma}\label{lem:lower-dimensional}
Let $(Y,B)$ be a projective klt $\Q$-pair with $1\le\dim Y<3$.
Let $T=K_Y+B$ be pseudo-effective and $M$ a nef rational Cartier divisor.
Assume the two lower-dimensional abundance hypotheses above.
If, for some $a\in\Q_{>0}$,
\[
 (T+aM)\cdot C>0
\]
for every integral curve through a sufficiently very general point of $Y$,
then $T+cM$ is big for every $c\in\Q_{>0}$.
\end{lemma}
\begin{proof}
Use the $M$-trivial adjoint reduction of
\cite[Section~8.2]{NS}, after clearing the Cartier index of $M$.
It gives a birational contraction to a klt pair $(Y',B')$ with
$T'=K_{Y'}+B'$ pseudo-effective, $M'$ nef, and
$A=T'+bM'$ nef for some rational $b>0$. On a common resolution,
\[
 p^*T=q^*T'+E,\qquad p^*M=q^*M',\qquad E\ge0
 \quad\text{$q$-exceptional}.
\]
The nef abundance hypothesis gives $A\equiv A_0$ with $A_0$ semiample.
If $A$ were not big, a positive-dimensional general fibre of the
morphism defined by $A_0$ would contain a curve $C$ through a very
general point, avoiding $q(\operatorname{Exc}q)$, with $A\cdot C=0$.
Pseudo-effectivity of $T'$ and nefness of $M'$ give
$T'\cdot C=M'\cdot C=0$. Its strict lift contradicts the positivity
of $T+aM$. Hence $A$ and $T+bM$ are big.
For $c>0$, choose $0<\lambda<\min\{1,c/b\}$ and write
\[
 T+cM=\lambda(T+bM)+(1-\lambda)T+(c-\lambda b)M.
\]
The first summand is big and the other two are pseudo-effective.
\end{proof}

\section{Intersection and base-locus estimates}

For an integral subvariety $V$ of a smooth variety and an effective
Cartier divisor $F$, $\mult_V(F)$ denotes the multiplicity at a
general point of $V$. If $f$ is a local equation at $\eta_V$, then
\[
 \mult_V(F)=\max\{j\ge0:f\in\mathfrak m_{\eta_V}^{j}\}.
\]
On a resolution of a subvariety, intersection products use the
pullbacks of the indicated restricted Cartier divisors.

\begin{lemma}\label{lem:hd-base-degree}
Let $X$ be a smooth projective variety of dimension $n$, and let
$D$ be a nef Cartier divisor. Let $\Lambda\subseteq|D|$ be a
nonempty linear system and $V$ an integral component of its base
locus, of codimension $c$. If $\mult_V(F)\ge h$ for every
$F\in\Lambda$, where $h$ is a positive integer, then
\[
 h^cD^{n-c}\cdot V\le D^n.
\]
\end{lemma}
\begin{proof}
Intersect successively with $c$ general members of $\Lambda$,
discarding the components contained in the base locus before
each further cut. The remaining components meet the next member
properly. Since $D$ is nef, discarding effective components does
not increase the degree measured by the remaining powers of $D$.
The resulting codimension-$c$ cycle therefore has $D$-degree at
most $D^n$.

At $\eta_V$, the base ideal is primary for the maximal ideal:
this uses that $V$ is a component of the base locus. No discarded
component contains $\eta_V$. The $c$ general local equations
$f_1,\ldots,f_c$ form a parameter ideal $I$ in the regular local
ring $R=\mathcal O_{X,\eta_V}$ of dimension $c$. Its local
intersection multiplicity is $\operatorname{length}(R/I)$.
The hypothesis gives $I\subseteq\mathfrak m^h$, so monotonicity
of Hilbert--Samuel multiplicity yields
\[
 \operatorname{length}(R/I)=e(I,R)
 \ge e(\mathfrak m^h,R)=h^c.
\]
Thus the coefficient of $V$ is at least $h^c$. Taking $D$-degree
and using nefness for all remaining effective components proves
the inequality.
\end{proof}

\begin{lemma}\label{lem:finite-subadjunction-model}
Let $W$ be smooth projective, $\Theta$ an effective rational divisor,
and $V\subset W$ an integral subvariety of positive dimension.
Suppose $(W,\Theta)$ is lc at $\eta_V$ and has an lc place centred on $V$.
Let $V^\nu\to V$ be the normalization and put
$L=(K_W+\Theta)|_{V^\nu}$. There are a finite surjective morphism
$\eta:V'\to V^\nu$, with $V'$ normal, an effective rational divisor
$B_{V'}$, and the trace $M_{V'}$ of a b-nef rational b-divisor such that
\[
 \eta^*L\sim_\Q K_{V'}+B_{V'}+M_{V'}.
\]
\end{lemma}
\begin{proof}
We use the construction in \cite[Lemma~7.1]{LA}.
The dlt blowup \cite[Theorem~2.10]{FH} gives
\[
 p^*(K_W+\Theta)=K_Q+B_Q+E_Q,
 \qquad (Q,B_Q)\text{ dlt},\quad E_Q\ge0,
\]
where $B_Q$ is the truncated boundary and the effective surplus
$E_Q$ is absent near $\eta_V$. Choose a dlt stratum $S$ minimal
among those dominating $V$. Adjunction, followed by adding
$E_Q|_S$, gives an effective rational divisor $B_S$ and
\[
 K_S+B_S\sim_\Q(p|_S)^*((K_W+\Theta)|_V).
\]
The stratum is not contained in $E_Q$. Minimality makes
$(S,B_S)$ klt over $\eta_V$, although it need not be globally lc.

Factor through the normalization and the connected-fibre base:
$S\xrightarrow{q}V'\xrightarrow{\eta}V^\nu\to V$, with $\eta$
finite. The discrepancy-round-up rank condition for $q$ is one:
over the generic point the effective klt boundary contributes
zero round-up on strict divisors and allows poles only along
exceptional divisors. Its pushforward to normal $S$ is therefore
$\OO_S$, and connectedness gives rank one. Thus this is a
klt-trivial fibration in \cite[Definition~3.1]{FG}. Its formula is
\[
 \eta^*((K_W+\Theta)|_{V^\nu})
 \sim_\Q K_{V'}+B_{V'}+M_{V'},
\]
with b-nef moduli divisor $M$ by \cite[Theorem~3.5]{FG}.
The discriminant on this initial base is effective: over a prime
divisor, an inverse-image component of multiplicity $m\ge1$ and
boundary coefficient $b\ge0$ bounds its lc threshold by
$(1-b)/m\le1$. A negative threshold still gives a positive
discriminant coefficient. Only generic sub-klt singularities are
required for this use of the formula.
\end{proof}

\begin{lemma}\label{lem:generic-subadjunction}
Let $W$ be smooth projective, $\Theta\ge0$ a rational divisor,
and $V\subset W$ an integral subvariety of dimension $f>0$.
Suppose $(W,\Theta)$ is lc at $\eta_V$ and has an lc place centred on $V$.
For every nef rational Cartier divisor $P_V$ on $V$ and resolution
$\rho:Y\to V$,
\[
 K_Y\cdot(\rho^*P_V)^{f-1}
 \le (K_W+\Theta)|_V\cdot P_V^{f-1}.
\]
\end{lemma}
\begin{proof}
Apply Lemma~\ref{lem:finite-subadjunction-model} and intersect with
$f-1$ copies of the pulled-back $P_V$. The effective discriminant
and the moduli trace have nonnegative intersection; for the latter,
push down from a higher base where the moduli divisor is nef.
For compatible canonical Weil cycles,
\[
 \eta_*K_{V'}=(\deg\eta)K_{V^\nu}+\eta_*R_\eta,
 \qquad R_\eta\ge0.
\]
Discard ramification and divide by $\deg\eta$. Pushing down $K_Y$
to $V^\nu$ changes it only by exceptional cycles, whose testing
intersection is zero. No rational Cartier assumption on
$K_{V^\nu}$ is needed.
\end{proof}

\begin{lemma}\label{lem:canonical-centre}
Let $W$ be smooth projective, $P_W$ a nef rational Cartier divisor,
and $k>0$ an integer with $kP_W$ Cartier. Let
$\Lambda\subseteq|kP_W|$ be a nonempty linear system and $V$ an
integral base-locus component of positive dimension and codimension $a$.
If $\mult_V(F)\ge h$ for every $F\in\Lambda$, where $h>0$ is an integer,
then there exists $c\in\Q_{>0}$ with $c\le ak/h$ such that, on a
resolution $\rho:Y\to V$,
\[
 K_Y\cdot(\rho^*(P_W|_V))^{\dim V-1}
 \le (K_W+cP_W)|_V\cdot(P_W|_V)^{\dim V-1}.
\]
\end{lemma}
\begin{proof}
The generic log canonical threshold of the base ideal satisfies
$\lambda\le a/h$, as seen by blowing up the smooth generic
centre. On a log resolution, average sufficiently many general
members of the subseries with total coefficient $\lambda$.
This gives an actual effective boundary
$\Theta\sim_{\Q}k\lambda P_W$ which is lc at the generic
point of $V$ and has an lc place centred on $V$.
Lemma~\ref{lem:generic-subadjunction} gives the inequality with
$c=k\lambda$. The construction uses actual general members of
the given subseries, and requires lc only at $\eta_V$.
Both the threshold bound and subadjunction are dimension-independent.
\end{proof}

\begin{lemma}\label{lem:jet-normal-containment}
Let $W$ be smooth projective, $L$ a line bundle, and $z:T\to W$
a morphism from a smooth integral variety. Fix integers
$m_1>m_0\ge0$ and $1\le h\le m_1-m_0+1$.
Suppose the spaces
\[
 \mathcal H_i(t)=H^0(W,L\otimes\mathfrak m_{z(t)}^{m_i})
 \quad(i=0,1)
\]
form vector bundles, and a reduced family $\Gamma\subset T\times W$
dominates $W$, with each general fibre $\Gamma_t$ a component of
both base loci. Then the higher base ideal satisfies
$\mathfrak b_{1,t}\subset I_{\Gamma_t}^{h}$ at a general incidence point.
\end{lemma}
\begin{proof}
A frame of the higher kernel has the form
$s(t,w)=\sum_\alpha a_\alpha(t)s_\alpha(w)$ for fixed global sections.
A parameter derivative of order $j<h$ loses at most $j$ orders
at the moving mark, so belongs to the full lower kernel and
vanishes on $\Gamma$.
At a general smooth point, dominance of $\Gamma\to W$ makes
parameter directions span the normal space. Choose transverse
parameters $t'$ and write $\Gamma$ locally as
$u=t'-F(t'',w)=0$. Differentiation in $t'$ with $t'',w$ fixed
is differentiation in $u$. Thus all terms of degree less than $h$
in the normal parameters vanish, giving $s\in I_\Gamma^h$.
Characteristic zero and faithful flatness of the analytic local
ring give the algebraic containment. Restricting the finite frame
to a general parameter fibre proves the assertion, as in
\cite[Lemma~7.3]{LA}.
\end{proof}

\begin{lemma}\label{lem:restricted-components}
Let $g:W\to B$ be a morphism with $W$ smooth projective, and let
$\Lambda\subseteq|L|$ be a nonempty linear system with base ideal
$\mathfrak b$. Let $V$ be an integral base-locus component and
suppose $\mathfrak b\subset I_V^h$ near $\eta_V$, with $h>0$ an integer.
For general $b\in g(V)$, let $F$ be a reduced component of $V\cap W_b$,
where $W_b=g^{-1}(b)$. If $W_b$ is smooth and integral and
$F\subsetneq W_b$, then the restricted system on $W_b$ is nonzero,
$F$ is a base-locus component, and its base ideal satisfies
\[
 \mathfrak b\mathcal O_{W_b}\subset I_F^h
 \quad\text{near }\eta_F.
\]
\end{lemma}
\begin{proof}
Remove the other base components and work where $V\to g(V)$ is
smooth over its actual image. Its scheme fibre is reduced there,
hence $I_V\OO_{W_b}=I_F$ and $I_V^h\OO_{W_b}=I_F^h$.
The local base support is exactly $F$.
Since $F$ is proper, not all restricted sections vanish identically
on $W_b$. The resulting base ideal is primary at $\eta_F$.
This is the actual-image restriction in \cite[Lemma~7.4]{LA}.
\end{proof}

\begin{lemma}\label{lem:full-kernels}
Let $W$ be a smooth projective variety of dimension $d$, and let
$P_W$ be ample and rational Cartier. Fix an integer $k>0$ with $kP_W$ Cartier and levels
$0\le\tau_0<\cdots<\tau_d$ with common gap $\delta>0$ and
$k\delta\ge4$. Suppose that, at a very general marked point $z$,
all full spaces
\[
 \mathcal H_i(z)=H^0(W,\OO_W(kP_W)\otimes\mathfrak m_z^{\lceil k\tau_i\rceil})
\]
are nonzero, and the lowest base locus has a positive-dimensional
component through $z$. Then there is an algebraic covering family
of proper integral components $V_t$, common to two adjacent base
loci, with $1\le f=\dim V_t<d$. At a general incidence point the
entire higher-kernel base ideal satisfies
\[
 \mathfrak b_t\subset I_{V_t}^h,\qquad
 h=\lceil k\delta/2\rceil,\qquad k/h\le2/\delta.
\]
Consequently, on a resolution $V_t^*$,
\begin{equation}\label{eq:full-kernel-bounds}
 \begin{split}
 0<P_W^f\cdot V_t&\le(2/\delta)^{d-f}P_W^d,\\
 K_{V_t^*}\cdot P_W^{f-1}&\le(K_W+cP_W)|_{V_t}\cdot P_W^{f-1},
 \qquad 0<c\le2(d-f)/\delta.
 \end{split}
\end{equation}
\end{lemma}
\begin{proof}
On the constant-rank open, the full kernels are vector bundles
with nested proper base loci. Track a component through the mark
across the $d+1$ levels. Its dimensions lie in
$\{1,\ldots,d-1\}$, so two adjacent tracked components coincide.
After a finite extension of the parameter field, spread them out
to an integral family whose marked incidence dominates $W$.
The adjacent orders satisfy
\[
 \lceil k\tau_{i+1}\rceil-\lceil k\tau_i\rceil
 \ge k\delta-1\ge h-1.
\]
Lemma~\ref{lem:jet-normal-containment} gives the ordinary-power
containment for the entire higher base ideal. Apply
Lemmas~\ref{lem:hd-base-degree} and \ref{lem:canonical-centre}
for the degree and canonical bounds. The construction follows
\cite[Lemmas~4.5--4.6]{NS}.
\end{proof}

\begin{lemma}\label{lem:bounded-curves}
For each $f\in\{1,2,3\}$, there is an integer $m_f>0$ depending
only on $f$ with the following property. Let $Y$ be smooth projective
of dimension $f$. Let $B_0$ be a reduced simple normal crossing
divisor or zero, $L$ a nef Cartier divisor of full nef dimension,
and $P_Y$ a nef rational Cartier divisor. Let $r,C_0,C_1>0$. If
\[
 0<P_Y^f\le C_0,\quad P_Y-rL\text{ is nef},\quad
 K_Y+B_0+L\text{ is big},
\]
and
\[
 (K_Y+B_0+L)\cdot P_Y^{f-1}\le C_1P_Y^f,
\]
then
\[
 r\le(m_fC_1)^{f-1}C_0,
\]
where $m_f$ depends only on $f$.
\end{lemma}
\begin{proof}
Birkar--Zhang \cite[Theorem~1.3]{BZ}, with boundary coefficients
in $\{0,1\}$ and Cartier index one for $L$, makes
$|m_f(K_Y+B_0+L)|$ birational. Resolve its base ideal and let
$H$ be the free moving part. Then $H$ is nef and big and
\[
 H\cdot P_Y^{f-1}\le m_fC_1P_Y^f.
\]
Khovanskii--Teissier gives
\[
 P_Y\cdot H^{f-1}\le P_Y^f
 \left(\frac{H\cdot P_Y^{f-1}}{P_Y^f}\right)^{f-1}
 \le(m_fC_1)^{f-1}C_0.
\]
General hyperplanes through the image of a very general point
under the birational morphism defined by $H$ cut an effective
curve cycle containing a curve $C$ through that point.
Full nef dimension and integrality give $L\cdot C\ge1$.
Thus $r\le P_Y\cdot C\le P_Y\cdot H^{f-1}$. For $f=1$ take $C=Y$.
This is \cite[Lemma~4.1]{NS}; the constants depend only on the dimension.
\end{proof}

\section{Proof of Lemma~\ref{lem:full-jets}}\label{sec:large-jets}



\subsection{The projective bundle and adjoint bigness}

In this section, write $D_i=(p_i\circ\pi)^*D$ and
$P_i=(p_i\circ\pi)^*P$ for the technical intersection calculations.
For $P=P_n$ and $s>0$, write $R=R_{n,s}$. On a slice over one fixed
base coordinate, put
\[
 Z_{\rm sl}=\mathbb P(\OO_X\oplus\OO_X(D)),\qquad
 R_{\rm sl}=P+s\xi.
\]
We use $\pi_*\OO_Z(a)=\operatorname{Sym}^a
(\OO_X(D)_1\oplus\OO_X(D)_2)$; thus $\dim Z=7$ and
$\dim Z_{\rm sl}=4$.

\begin{lemma}\label{lem:bundle-volume}
With $P_a=P+aD$, the projective-bundle formulas give
\[
 R^7=140\int_0^s P_a^3P_{s-a}^3\,da,\qquad
 R_{\rm sl}^4=4\int_0^s P_a^3\,da.
\]
For fixed $s$, these tend to $140s\epsilon^6$ and
$4s\epsilon^3$, respectively, as $n\to\infty$. Also
\[
 K_Z\sim D_1+D_2-2\xi,\qquad K_{Z_{\rm sl}}\sim D-2\xi.
\]
\end{lemma}
\begin{proof}
For a split rank-two bundle,
$\pi_*(\xi^j)=\sum_{a+b=j-1}D_1^a\cdot D_2^b$ for $j\ge1$.
Expand $R^7$ and integrate the binomial terms to obtain the first
identity; its coefficient is $7!/(3!3!)=140$.
The analogous one-factor expansion gives the second identity.
Since $D^3=0$, $P^3\to\epsilon^3$, and $r_n\to\infty$,
$(P+aD)^3/P^3\to1$ uniformly for $0\le a\le s$.
The two limits follow. The canonical formulas are the rank-two
projective-bundle formula with $K_X\sim0$.
\end{proof}

\begin{lemma}\label{lem:semiample-trivial-rays}
Let $U,M'$ be rational Cartier divisors on a normal projective
variety, with $U$ semiample and $dU$ Cartier and free for an
integer $d>0$. Let $0<t_0<t$ satisfy
\[
 (1-t_0/t)/d>12t_0/t.
\]
If $C$ generates a $(U+t_0M')$-negative extremal ray and
$(U+tM')\cdot C\ge-12$, then $U\cdot C=0$. A contraction or flip of such
a ray preserves the semiampleness of $U$.
\end{lemma}
\begin{proof}
If $U\cdot C>0$, integrality gives $U\cdot C\ge1/d$, and
\[
 (U+t_0M')\cdot C=(1-t_0/t)U\cdot C+(t_0/t)(U+tM')\cdot C
 \ge(1-t_0/t)/d-12t_0/t>0,
\]
a contradiction. Thus $U\cdot C=0$ on every contracted ray.
The free line bundle $dU$ descends through a contraction and
pulls back through a flip, preserving semiampleness.
\end{proof}

\begin{lemma}\label{lem:positive-boundary-model}
Let $X$ be a smooth projective threefold with $K_X\sim0$,
and let $D$ be a nef Cartier divisor with $\nef(X,D)=3$.
Let $G\ge0$ be rational with $(X,G)$ klt and $\kappa(G)>0$.
There exist $t_0>0$ rational, a normal projective model $Y$, and
a common resolution $p:W\to X$, $q:W\to Y$, with
$U=K_Y+G_Y$ semiample, $Q_0=U+t_0M'$ nef, and
\begin{equation}\label{eq:positive-boundary-comparison}
 p^*G=q^*U+E_0,\qquad
 p^*(G+t_0D)=q^*Q_0+E_1,
 \qquad E_0,E_1\ge0\text{ $q$-exceptional}.
\end{equation}
Moreover, $\kappa(U)=\kappa(G)$, $\nef(Y,Q_0)=3$, and $M'$
is the transform of $D$, viewed as the trace of its nef b-divisor.
\end{lemma}
\begin{proof}
We spell out \cite[Proposition~3.3]{NS} in the only threefold
instance required below. Run the ordinary $G$-MMP to a good
model; ordinary threefold abundance is \cite{KollarEtAl}.
Its adjoint $U=K_{Y_0}+G_{Y_0}$ is semiample. Let $M'_0$ be the
trace of the fixed nef b-divisor determined by $D$ on $X$.
For small $t>0$, this finite first program is also negative for
$G+tD$, and $(Y_0,G_{Y_0}+t\mathbf D)$ is generalized klt.


Choose $dU$ free and $t_0$ as in
Lemma~\ref{lem:semiample-trivial-rays}. Run a generalized
$(U+t_0M'_0)$-program. The nef data are rational and NQC;
ordinary klt approximation gives its steps, and
\cite[Theorem~1.2]{ChenTsakanikas} gives termination.
The generalized length bound is $12$ in this threefold setting
\cite[Lemma~3.1]{NS}. Lemma~\ref{lem:semiample-trivial-rays}
makes every step $U$-trivial.
On the final model $Y$, $U$ remains
semiample, $Q_0=U+t_0M'$ is nef, and a common resolution gives the two displayed comparisons.

Moreover $Q_0$ has full nef dimension. Otherwise a very general
compact nef-reduction fibre has a curve through its very general
point avoiding $q(\operatorname{Exc}q)$, which has codimension
at least two. Nefness of $U$, pseudo-effectivity of $M'$, and
generic nonnegativity of pseudo-effective degrees give $U\cdot C=M'\cdot C=0$.
Lifting this curve in both comparisons gives $D$-degree zero
through a very general point of $X$, a contradiction.
\end{proof}

\begin{lemma}\label{lem:positive-boundary}
Let $X$ be a smooth projective threefold with $K_X\sim0$,
and let $D$ be a nef Cartier divisor with $\nef(X,D)=3$.
If $G\ge0$ is rational, $(X,G)$ is klt, and $\kappa(G)>0$,
then $G+cD$ is big for every rational $c>0$.
\end{lemma}
\begin{proof}
Use the model in Lemma~\ref{lem:positive-boundary-model}.
The contraction $h:Y\to S$ defined by $U$ has
$\dim S=\kappa(G)>0$ and $U\sim_\Q h^*A_S$ with $A_S$ ample.
On a very general positive-dimensional fibre $F$,
\[
 K_F+G_Y|_F\sim_\Q0,
 \qquad Q_0|_F=t_0M'|_F\text{ is nef of full nef dimension}.
\]
Its dimension is at most two, so Lemma~\ref{lem:lower-dimensional}
makes $Q_0|_F$ big. For zero-dimensional fibres relative bigness
is automatic. Relative bigness follows on the generic fibre: for a fixed ample
$H$, one of the countably many closed section loci for
$(\ell Q_0-H)|_{F_t}$ contains a very general point, hence is the
whole shrunken flat base. Adding a sufficiently large multiple
$bU$ now makes $Q_0+bU$ big globally. The comparisons imply
that $(1+b)G+t_0D$ is big. Finally, for
$0<\lambda<\min\{1/(1+b),c/t_0\}$,
\[
 G+cD=\lambda((1+b)G+t_0D)
 +(1-\lambda(1+b))G+(c-\lambda t_0)D
\]
is big.
\end{proof}

\begin{lemma}\label{lem:multisection-axes}
Let $X$ be a smooth projective threefold with $K_X\sim0$, and let
$D$ be nef Cartier with $\nef(X,D)=3$. Let $P$ be ample and rational
Cartier, and let $s>0$ be an integer. Put $R_{\rm sl}=P+s\xi$ on the
quotient projective bundle below. Let $F\subset
\mathbb P(\OO_X\oplus\OO_X(D))$ be an integral multisection
of degree $e>0$ meeting the torus open. Writing $p$ for the bundle projection and $S_1,S_2$ for the summand
sections, set $G_i=p_*(S_i\cdot F)$. These are effective integral divisors $G_1,G_2$ on $X$ satisfying
\begin{equation}\label{eq:axis-divisors}
 G_1-G_2\sim\pm eD,\qquad
 sG_i\cdot P^2\le R_{\rm sl}^3\cdot F.
\end{equation}
For $G=G_1+G_2$ and every sufficiently small rational $b>0$,
$(X,bG)$ is klt, $\kappa(bG)>0$, and $bG+D$ is big.
\end{lemma}
\begin{proof}
The norm of the ratio of the two fibre coordinates gives the
linear equivalence. If $S$ is either axis, $R_{\rm sl}-sS$
and $R_{\rm sl}-P$ are nef. Intersecting these classes against
the effective cycle $S\cdot F$ gives the degree inequalities.

If $\kappa(bG)=0$, ordinary threefold abundance gives a good
$(K_X+bG)$-model $Y$ with $K_Y+bG_Y\sim_\Q0$.
The transform of $K_X$ is linearly trivial, so $G_Y=0$:
every component of $G$ has been contracted. On a common
resolution $a:U\to X$, $c:U\to Y$, the signed relation in
\eqref{eq:axis-divisors} expresses $a^*D\sim_\Q E$, with
$E$ $c$-exceptional. Divisors exceptional over $X$ are also
exceptional over $Y$, since this MMP extracts none.
Since $a^*D$ is nef, negativity gives $E\le0$.
Its class is pseudo-effective, so intersection with an ample
square forces $E=0$. This contradicts full nef dimension.
Thus $\kappa(bG)>0$, and Lemma~\ref{lem:positive-boundary}
gives bigness. The axis construction and negativity argument
are dimension-independent; only the ordinary good-model input
is being specialized to threefolds.
\end{proof}

\begin{lemma}\label{lem:fixed-graphs}
Let $T_1,T_2$ be non-uniruled integral projective varieties
over $\C$ of the same dimension. If an integral algebraic family
of graphs of birational maps $T_1\dashrightarrow T_2$ dominates
$T_1\times T_2$, both varieties are birational to an abelian variety.
\end{lemma}
\begin{proof}
Identify the two varieties by one map in the family.
Hanamura's theorem, in the form
\cite[Theorem~3.8]{Blanc}, supplies a smooth projective birational
model $Y$ with $\operatorname{Bir}(Y)_{\rm red}$ a group scheme
locally of finite type. Its identity component is
$\operatorname{Aut}^0(Y)$, an abelian variety acting regularly.

Conjugate the family to $Y$, take the closure of its relative
graphs, and shrink the integral parameter space so that these
graphs are flat and geometrically integral, with both projections
birational. This shrinking is necessary: graph flatness is not
preserved automatically by birational conjugation
\cite[Remark~3.6]{Blanc}. The flat-graph representability theorem
\cite[Proposition~3.7]{Blanc} then gives a morphism from the
reduced parameter variety to $\operatorname{Bir}(Y)_{\rm red}$.
Its connected image lies in one translate of the identity component.
After translating by a fixed member, product dominance implies
that a general point has a dense orbit under that abelian group.
Such an orbit is the quotient by its closed stabilizer, hence an
abelian variety. It is therefore a birational model of $Y$.
\end{proof}

\begin{lemma}\label{lem:abelian-adjoint}
Let $p:Y\to A_0$ be a birational morphism from a smooth projective
variety to an abelian variety. If $L$ is nef Cartier of full nef
dimension on $Y$, that is, $\nef(Y,L)=\dim Y$, then $K_Y+L$ is big.
\end{lemma}
\begin{proof}
The numerical pushforward $M'=p_*L$ is pseudo-effective on $A_0$,
thus nef. If it were not big, the kernel of the semipositive
Hermitian form of its integral class would define a
positive-dimensional abelian subvariety on whose translates $M'$
is numerically trivial. The image of the exceptional locus of $p$
has codimension at least two. A general translate meets that image
in codimension at least two; complete-intersection curves in the
translate can be chosen through a very general point avoiding it.
Their strict lifts have $L$-degree zero, a contradiction. Thus $M'$
is big.

Negativity gives $p^*M'-L=E_0\ge0$, with $E_0$ exceptional.
Since $A_0$ is smooth and $K_{A_0}\sim0$, $K_Y\sim E_K\ge0$,
with positive coefficient along every $p$-exceptional divisor.
Choose $0<u<1$ rational with $E_K-uE_0\ge0$. Then
\[
 K_Y+uL\sim_\Q up^*M'+(E_K-uE_0)
\]
is big. Adding $(1-u)L$ proves the assertion. This lemma has no
dimension restriction.
\end{proof}

\subsection{Base components and slice exclusions}

Fix $\delta=8\epsilon$, $\tau_i=9+i\delta$ for $0\le i\le7$,
and choose $s$ so that $\tau_7^7<70s\epsilon^6$.
For all sufficiently large $n$, Lemma~\ref{lem:bundle-volume} gives
\[
 \epsilon^3/2<P^3<2\epsilon^3,\qquad
 \tau_7^7<R^7<280s\epsilon^6,\qquad R_{\rm sl}^4<8s\epsilon^3.
\]
In particular $(2/\delta)^3\,8\epsilon^3=1/8<1$.
For the following lemmas, these choices remain fixed, and
$C_0=(2/\delta)^4\,280s\epsilon^6$.

\begin{lemma}\label{lem:sevenfold-centres}
For all sufficiently large $n$, if $\varepsilon(R;z)\le8$ at every
point of the torus open, there is a covering family of proper
integral $V\subset Z$, with $f=\dim V$, such that
\begin{equation}\label{eq:sevenfold-centres}
 0<R^f\cdot V\le(2/\delta)^{7-f}R^7,\qquad
 K_{V^*}\cdot R^{f-1}\le\frac{2(7-f)}{\delta}R^f\cdot V.
\end{equation}
For $i=1,2$, set $\pi_i=p_i\circ\pi$. For a general reduced fibre
component $F$ of $V\to\pi_i(V)$, viewed in the opposite fourfold slice, if
$1\le d_F=\dim F\le3$, the actual restricted system gives
\begin{equation}\label{eq:slice-centres}
 \begin{split}
 0<R_{\rm sl}^{d_F}\cdot F&\le(2/\delta)^{4-d_F}R_{\rm sl}^4,\\
 K_{F^*}\cdot R_{\rm sl}^{d_F-1}
 &\le(K_{Z_{\rm sl}}+c_{d_F}R_{\rm sl})|_F \cdot R_{\rm sl}^{d_F-1},
 \quad 0<c_{d_F}\le2(4-d_F)/\delta.
 \end{split}
\end{equation}
Both images of a general marked incidence point are very general on $X$.
\end{lemma}
\begin{proof}
Asymptotic Riemann--Roch and $\tau_7^7<R^7$ make the eight full
kernels nonzero for large divisible $k$. A curve through a very
general $z$ with $R$-degree divided by multiplicity less than
$\tau_0$ is contained in the lowest base locus, by the intersection
multiplicity inequality. Hence that base locus is positive-dimensional.
Lemma~\ref{lem:full-kernels} gives the covering family and the first
bounds. Since $K_Z$ is minus the axes and $V$ meets their complement,
the ambient canonical contribution is nonpositive.
Apply Lemma~\ref{lem:restricted-components} to either projection
over its actual image, and then apply
Lemmas~\ref{lem:hd-base-degree} and \ref{lem:canonical-centre}
in the smooth fourfold slice. This gives the fibre bounds.
\end{proof}

For the next three lemmas, $F$ is a general reduced fibre component
of $V\to\pi_i(V)$ in a covering family as in
Lemma~\ref{lem:sevenfold-centres}; $W_0$ is its image in the other copy
of $X$. We require $1\le\dim F\le3$, the slice estimates
\eqref{eq:slice-centres}, and the stated very-general-image conditions.

\begin{lemma}\label{lem:slice-small-image}
For all sufficiently large $n$, a slice component $F$ satisfying
the above conditions cannot be generically finite onto $W_0$ with $\dim W_0<3$.
\end{lemma}
\begin{proof}
If $F\to W_0$ is generically finite and $\dim W_0<3$, the
pullback $L$ of $D$ is nef Cartier of full nef dimension on $F^*$,
and $K_{F^*}$ is pseudo-effective by the very-general-images
argument. Lemma~\ref{lem:lower-dimensional} gives $K_{F^*}+L$
big. The testing class $P_Y=R_{\rm sl}|_{F^*}$ satisfies
$P_Y-r_nL$ nef, and \eqref{eq:slice-centres} gives uniform volume
and canonical ratios. Lemma~\ref{lem:bounded-curves} bounds $r_n$,
a contradiction.
\end{proof}

\begin{lemma}\label{lem:saturated-slices}
For all sufficiently large $n$, a slice component $F$ satisfying
the above conditions cannot be the full projective-line bundle over $W_0$ with $\dim W_0\le2$.
\end{lemma}
\begin{proof}
If $F\to W_0$ has one-dimensional fibres, $F$ is the entire
projective-line bundle over $W_0$. Put $w=\dim W_0=d_F-1$.
For $w=0$, \eqref{eq:slice-centres} gives
\[
 s=R_{\rm sl}\cdot F\le(2/\delta)^3R_{\rm sl}^4<s.
\]
For $w=1,2$, resolve $W_0$ and take its induced bundle
$p_0:F^*\to W_0^*$. Cancelling the relative canonical term
$-2\xi+p_0^*D$, which equals the ambient canonical restriction,
gives
\[
 p_0^*K_{W_0^*}\cdot R_{\rm sl}^w\le c_{d_F}R_{\rm sl}^{w+1}\cdot F.
\]
Its term linear in $\xi$ is $wsK_{W_0^*}\cdot P^{w-1}$; all other
terms are nonnegative because $K_{W_0^*}$ is pseudo-effective
and $P,D$ are nef. For $r_n\ge s$, expansion gives
\begin{equation}\label{eq:saturated-estimate}
 \begin{split}
 (w+1)sP^w\cdot W_0&\le R_{\rm sl}^{w+1}\cdot F\le7sP^w\cdot W_0,\\
 K_{W_0^*}\cdot P^{w-1}&\le(7c_{d_F}/w)P^w\cdot W_0.
 \end{split}
\end{equation}
The first lower bound and \eqref{eq:slice-centres} bound $P^w\cdot W_0$
uniformly. Applying Lemmas~\ref{lem:lower-dimensional} and
\ref{lem:bounded-curves} on $W_0^*$ is again contradictory.
\end{proof}

\begin{lemma}\label{lem:multisection-slices}
For all sufficiently large $n$, a slice component $F$ satisfying
the above conditions cannot be a multisection over $X$.
\end{lemma}
\begin{proof}
Lemma~\ref{lem:multisection-axes} gives $G=G_1+G_2$, $bG+D$ big,
and $G\cdot P^2\le2R_{\rm sl}^3\cdot F/s$. On a log resolution $a:Y\to X$,
let $B_0$ be the reduced strict support of $G$ together with every
exceptional divisor. Decrease $b$ so that
\[
 K_Y+B_0+a^*D\sim_\Q a^*(bG+D)+E_b,\qquad E_b\ge0.
\]
This adjoint is big. Testing with $P_Y=a^*P$, all exceptional
divisors contribute zero, and
\[
 (K_Y+B_0+a^*D)\cdot P_Y^2\le G\cdot P^2+D\cdot P^2\le C,
 \qquad D\cdot P^2\le P^3/r_n.
\]
The positive lower bound on $P^3$ gives a uniform adjoint ratio.
Lemma~\ref{lem:bounded-curves}, with this reduced boundary,
excludes the case.
\end{proof}

\begin{lemma}\label{lem:finite-correspondences}
For all sufficiently large $n$, both projections of a general
$V$ in Lemma~\ref{lem:sevenfold-centres} are generically finite
and dominant. Let $\rho:Y\to V$ be a smooth resolution, put
$g_i=\pi_i\circ\rho:Y\to X$ and $P_Y=\rho^*(R|_V)$. Then
\begin{equation}\label{eq:correspondence-estimates}
 0<P_Y^3\le C_0,\qquad K_Y\cdot P_Y^2\le(8/\delta)P_Y^3,
 \qquad P_Y-r_ng_i^*D\text{ is nef}.
\end{equation}
Moreover, $\deg g_i\le2C_0/\epsilon^3$.
\end{lemma}
\begin{proof}
Lemmas~\ref{lem:slice-small-image}, \ref{lem:saturated-slices},
and \ref{lem:multisection-slices} exclude fibre dimensions one,
two, and three. A four-dimensional fibre makes $V$ the full
inverse image of a proper $W_0\subset X$, with $\dim W_0\le2$;
the other projection then has fibre dimension $\dim W_0+1$, already
excluded. Thus both projections are generically finite and $f\le3$.
For $f<3$, \eqref{eq:sevenfold-centres} and
Lemmas~\ref{lem:lower-dimensional} and \ref{lem:bounded-curves}
give a uniform bound on $r_n$, a contradiction. Hence $f=3$.
The displayed estimates follow from \eqref{eq:sevenfold-centres}
and nefness of $R-r_nD_i$.
Since $R-P_i$ is nef, $\deg g_i\,P^3\le P_Y^3\le C_0$.
\end{proof}


\subsection{Finite covers and the final exclusion}

\begin{lemma}\label{lem:branch-surfaces}
For all sufficiently large $n$, no relative Jacobian component
in the family of Lemma~\ref{lem:finite-correspondences} whose
general image is a divisor has images sweeping $X$.
\end{lemma}
\begin{proof}
Suppose that the images of one branch-divisor component sweep
$X$, and let $J\subset X$ be a general such image. Pulling back
a nowhere-vanishing three-form from $X$ gives an effective
canonical divisor on $Y$, containing a ramification divisor
over $J$ with positive integral coefficient. Thus
\[
 0<d_J:=P^2\cdot J\le K_Y\cdot P_Y^2\le(8/\delta)C_0.
\]
Generic numerical subadjunction for $(X,J)$ and Hodge index give
on a resolution $J^*$
\begin{equation}\label{eq:branch-surface}
 K_{J^*}\cdot P\le J^2\cdot P
 \le\frac{(J\cdot P^2)^2}{P^3}\le C_2d_J,
 \qquad C_2=16C_0/(\delta\epsilon^3).
\end{equation}
Here $(X,J)$ is lc at $\eta_J$; global lc is unnecessary.
Hodge index applies to $J$ without assuming $J$ nef.
Since the images sweep, $J^*$ is non-uniruled and the pulled-back
$D$ has full nef dimension. The testing volume is $d_J$, bounded
above, its canonical ratio is bounded by $C_2$, and $P-r_nD$ is
nef. Lemmas~\ref{lem:lower-dimensional} and
\ref{lem:bounded-curves} give a uniform bound for $r_n$.
Hence sweeping branch divisors are impossible for large $n$.
\end{proof}

\begin{lemma}\label{lem:finite-stein-covers}
Let $X$ be a smooth projective complex variety, $U\subset X$ a
nonempty Zariski open subset, and $e>0$ an integer. Up to isomorphism
over $X$, there are only finitely many finite surjective morphisms
$h:T\to X$ with $T$ normal and integral, $\deg h\le e$, and
$h^{-1}(U)\to U$ \'etale.
\end{lemma}
\begin{proof}
The fundamental group of $U$ is finitely generated, hence has
only finitely many subgroups of index at most $e$. These classify
the connected finite \'etale covers of $U$. Each normal cover of
$X$ is uniquely determined by its function field, as the normalization
of $X$ in that field. The assertion follows; compare
\cite[Lemma~7.5]{LA}.
\end{proof}

For a fixed family in Lemma~\ref{lem:finite-correspondences},
Lemma~\ref{lem:branch-surfaces} puts the finitely many relative
divisorial branch images in a proper closed subset of $X$.
Remove this subset and shrink the parameter space. Purity makes
each finite normal Stein cover \'etale over the resulting fixed open
$U$. The degree bound and Lemma~\ref{lem:finite-stein-covers} therefore
give finite lists of Stein covers for both projections.

\begin{lemma}\label{lem:bounded-correspondence-exclusion}
Let $X$ be a smooth projective threefold with $K_X\sim0$, and let
$D$ be a nef Cartier divisor with $\nef(X,D)=3$.
Consider an algebraic family of smooth projective threefolds $Y_t$
with generically finite morphisms $g_{i,t}:Y_t\to X$, for $i=1,2$,
whose product images cover a dense subset of $X\times X$.
Suppose that the normal Stein covers of each $g_{i,t}$ belong to
a finite list. Let $r,C_0,C_1>0$. If there are nef rational Cartier
divisors $P_t$ such that, for every general $t$,
\[
 0<P_t^3\le C_0,\qquad K_{Y_t}\cdot P_t^2\le C_1P_t^3,\qquad
 P_t-rg_{i,t}^*D\text{ is nef}\quad(i=1,2),
\]
then
\[
 r\le\max\{1,m_3^2(C_1+1)^2C_0\},
\]
where $m_3$ is the integer in Lemma~\ref{lem:bounded-curves}.
\end{lemma}
\begin{proof}
Each $Y$ is birational to both of its Stein covers. Fix a pair
$T_1,T_2$ occurring among the original members satisfying
the displayed bounds. The projected incidence
dominates $X\times X$. The possible graph families for the
finitely many pairs belong to countably many Hilbert schemes;
stratify for flat integral fibres and birational projections.
Uncountability of $\C$ implies that one such integral graph
family has evaluation dominant over $X\times X$. Its evaluation
is also dominant on $T_1\times T_2$: this product is integral,
finite over $X\times X$, and of the same dimension.

The covers are non-uniruled, since they are finite over $X$.
Lemma~\ref{lem:fixed-graphs} makes them birational to an abelian
variety $A_0$. Return to an \emph{original} correspondence for
this pair and resolve further to get a birational morphism
$p:Y\to A_0$. This preserves the displayed bounds; additional exceptional
canonical divisors have zero intersection with the pulled-back
testing square. The divisor $L=g_1^*D$ remains nef Cartier of
full nef dimension. Lemma~\ref{lem:abelian-adjoint} gives
$K_Y+L$ big, while
\[
 (K_Y+L)\cdot P_Y^2\le(C_1+1/r)P_Y^3.
\]
If $r<1$, the asserted bound is immediate. Otherwise
Lemma~\ref{lem:bounded-curves}, with $C_1+1/r\le C_1+1$, gives
$r\le m_3^2(C_1+1)^2C_0$. The sizes of the cover lists do not enter
this numerical bound.
\end{proof}


\begin{proposition}\label{prop:sevenfold-jets}
There is a fixed integer $s>0$ such that, for all sufficiently
large $n$, some point $z$ in the torus open satisfies
$\varepsilon(R;z)>8$.
\end{proposition}
\begin{proof}
Use the fixed choices above. If the assertion failed for arbitrarily
large $n$, Lemmas~\ref{lem:sevenfold-centres} and
\ref{lem:finite-correspondences} would produce bounded threefold
correspondences. Resolve and shrink the family, and spread out
the finitely many generic Jacobian components after a finite
parameter extension. Lemmas~\ref{lem:branch-surfaces} and
\ref{lem:finite-stein-covers} give finite lists of Stein covers.
Apply Lemma~\ref{lem:bounded-correspondence-exclusion} with
$r=r_n$ and $C_1=8/\delta$. Its bound contradicts $r_n\to\infty$.
\end{proof}


\begin{proof}[Proof of Lemma~\ref{lem:full-jets}]
Choose $z$ as in Proposition~\ref{prop:sevenfold-jets}, and rational
$8<a<\varepsilon(R;z)$. On the blowup $b:\widetilde Z\to Z$,
$b^*R-aJ_z$ is ample. For large divisible $k_0$, Serre vanishing
makes restriction surjective modulo $\mathfrak m_z^{k_0a}$,
with $k_0a\ge8k_0+1$. Thus $|k_0R_{n,s}|$ separates $8k_0$-jets at $z$.
\end{proof}

\begin{thebibliography}{99}
\bibitem{LA}
\emph{Log abundance in characteristic zero},
OpenAI mathematics collection, September 24, 2026,
Lemmas~7.1--7.5 (Tools for moving base components).
\href{https://github.com/openai/math/blob/main/preprints/Log-abundance-in-characteristic-zero-September-24-2026/paper.pdf}{Source manuscript}.
\bibitem{FH}
O.~Fujino and K.~Hashizume,
\emph{Existence of log canonical modifications and its applications},
Eur. J. Math. \textbf{9} (2023), Paper No.~13, Theorem~2.10.
\href{https://doi.org/10.1007/s40879-023-00598-0}{DOI}.
\bibitem{NS}
\emph{Numerical Semiampleness of Nef Adjoint Divisors},
OpenAI mathematics collection, October 3, 2026,
Proposition~3.3, Lemmas~4.1--4.7, Proposition~5.3,
Theorem~6.1, and Sections~8.4--8.5.
\href{https://github.com/openai/math/blob/main/preprints/Numerical-Semiampleness-of-Nef-Adjoint-Divisors-October-3-2026/paper.pdf}{Source manuscript}.
\bibitem{LP}
V.~Lazi\'c and Th.~Peternell,
\emph{On Generalised Abundance, I},
Publ. Res. Inst. Math. Sci. \textbf{56} (2020), 353--389.
Numbering follows arXiv:1808.00438v2, especially Corollary~C
and Lemmas~2.14 and 3.1.
\href{https://arxiv.org/pdf/1808.00438}{Preprint}.
\bibitem{FG}
O.~Fujino and Y.~Gongyo,
\emph{On the moduli $b$-divisors of lc-trivial fibrations},
Ann. Inst. Fourier \textbf{64} (2014), 1721--1735,
Definition~3.1 and Theorem~3.5.
\href{https://doi.org/10.5802/aif.2894}{DOI}.
\bibitem{BZ}
C.~Birkar and D.-Q.~Zhang,
\emph{Effectivity of Iitaka fibrations and pluricanonical
systems of polarized pairs},
Publ. Math. Inst. Hautes \'Etudes Sci. \textbf{123} (2016),
283--331, Theorem~1.3.
\href{https://doi.org/10.1007/s10240-016-0080-x}{DOI}.
\bibitem{KollarEtAl}
J.~Koll\'ar et al.,
\emph{Flips and Abundance for Algebraic Threefolds},
Ast\'erisque \textbf{211} (1992).
\href{https://www.numdam.org/item/AST_1992__211__1_0/}{Volume}.
\bibitem{ChenTsakanikas}
G.~Chen and N.~Tsakanikas,
\emph{On the termination of flips for log canonical generalized pairs},
Acta Math. Sin. (Engl. Ser.) \textbf{39} (2023), 967--994,
Theorem~1.2.
\href{https://arxiv.org/pdf/2011.02236}{Preprint}.
\bibitem{Blanc}
J.~Blanc,
\emph{Algebraic structures of groups of birational transformations},
Proc. Sympos. Pure Math. \textbf{94} (2017), 17--30,
Theorems~3.8--3.9.
\href{https://algebra.dmi.unibas.ch/blanc/articles/structurebir.pdf}{Published offprint}.
\end{thebibliography}

\end{document}
